\documentclass[12pt]{article}
\usepackage{float}
\usepackage{booktabs}
\usepackage{amsmath}
\usepackage{amssymb}
\usepackage[thmmarks, amsmath, thref]{ntheorem}
\usepackage{xcolor}
\usepackage{amsfonts}
\usepackage{geometry}
%\geometry{left=2.0cm, right=2.0cm, top=2.0cm, bottom=2.0cm}

\newtheorem{theorem}{Theorem}
\newtheorem{definition}{Definition}
\newtheorem{problem}{Problem}
\newtheorem{lemma}{Lemma}
\theorembodyfont{\color{blue}}
\theoremsymbol{\ensuremath{\blacksquare}}
\newtheorem{solution}{Solution}

\newcommand{\mat}[1]{\left(\begin{matrix}#1\end{matrix}\right)}
\newcommand{\arr}[2][n]{#2_1,#2_2,\cdots,#2_{#1}}
\newcommand{\iprod}[2]{\langle#1,#2\rangle}
\newcommand{\lr}[1]{\langle#1\rangle}
\newcommand{\mbb}[1]{\mathbb{#1}}
\newcommand{\mbf}[1]{\mathbf{#1}}
\newcommand{\mcal}[1]{\mathcal{#1}}
\newcommand{\mrm}[1]{\mathrm{#1}}
\newcommand{\on}[1]{\operatorname{#1}}
\newcommand{\ol}[1]{\overline{#1}}
\newcommand{\red}[1]{\textcolor{red}{#1}}
\newcommand{\blue}[1]{\textcolor{blue}{#1}}
\newcommand{\green}[1]{\textcolor{green}{#1}}

\DeclareMathOperator{\End}{End}
\DeclareMathOperator{\Ker}{Ker}
\DeclareMathOperator{\Img}{Im}
\DeclareMathOperator{\Hom}{Hom}
\DeclareMathOperator{\Tr}{Tr}
\DeclareMathOperator{\re}{Re}
\DeclareMathOperator{\im}{Im}
\DeclareMathOperator{\Sym}{Sym}
\DeclareMathOperator{\Rat}{Rat}
\DeclareMathOperator{\Rnk}{Rank}

\title{note template}
\date{\today}
\author{Cong Wen}

\begin{document}
\maketitle
\abstract{
This is another note calculating intersection problems involved in enumerative geometry, and Chern class is utilized in computation.
}

\section{Problem}
In projective space $\mbb{P}^2$(assuming the field is $\mbb{C}$), how many general conics are tangent to five general line?

We translate the problem into parameter space, every conic in the form of $a_0x_0^2+a_1x_1^2+a_2x_2^2+a_3x_0x_1+a_4x_0x_2+a_5x_1x_2$ corresponds bijectively to a point in $X:=\mbb{P}^5$, although there are some of them we want to rule out.

For each of the five general lines, the condition of tangency corresponds to a quadratic surface $Z(Q_i)$ in $X$, in which $Q_i$ is a quadratic polynomial in terms of homogeneous coordinate. Then we have to count points in

$$Z:=\bigcap\limits_{i=1}^5Z(Q_i)\subset X.$$

Using the rational equivalence in $\mbb{P}^n$, the Chow ring gives $A^*(\mbb{P}^n)=\mbb{Z}[H]/H^{n+1}$. In this case, $[Z(Q_i)]=2H$, in which $H$ is the hyperplane class of $\mbb{P}^5$. Now we get the number of points of $Z$
$$\#(Z)=\int_{X}(2H)^5=32$$

But there are some cases we need to rule out, $Z'$ as a subset of $Z$ represents double lines which trivially satisfy the tangency conditions, now we need to compute its contribution and subtract it from $32$.


In fact, since $Q_i$ is quadratic and expressed in terms of homogeneous coordinates, it can be viewed as a section of $\mcal{O}_{X}(2)$, then $Z$ is as same as the zero locus of $s=(s_1,\cdots,s_5)\in\mcal{O}_{X}(2)^{\oplus 5}$, this also gives that $c(\mcal{O}_{X}(2)^{\oplus 5})=(1+2H)^5$, and

$$\#Z=\int_{X}c_{\rm{top}}(\mcal{O}_{X}(2)^{\oplus 5})=32.$$

Now consider the subspace $V\subset X$ which represents double lines, for points $(a_0,\cdots,a_5)$ in $V$, $a_0x_0^2+a_1x_1^2+a_2x_2^2+a_3x_0x_1+a_4x_0x_2+a_5x_1x_2=0$ must be reduced to $(b_0x_0+b_1x_1+b_2x_2)^2=0$, which is in fact isomorphic to $\mbb{P}^2$, so we have $i:\mbb{P}^2\cong V\hookrightarrow X$ in which $V$ is a quadratic surface in $X$, so the pullback $i^*:A^*(X)\rightarrow A^*(V)$ gives rise to $i^*(H)=2h$, $h$ is the hyperplane class in $V$.

Then we need to compute number of $Z(s|_V)=Z'\subset V$ mentioned above, to do this, we identify $s|_V$ as a general section of line bundles on $V$ then make use of Chern class computation.

By adjusting coordinate system locally, $TX=\lr{\partial_{y_0},\partial_{y_1},\partial_{y_2},\partial_{y_3},\partial_{y_4},\partial_{y_5}}$ and $TV=\lr{\partial_{y_0},\partial_{y_1},\partial_{y_2}}$, in which $(y_0,y_1,y_2)$ is the local coordinate of $V$.

Computation Part


We need a lemma,
\lemma{$U$ is a open neighborhood of $0$ in $\mbb{C}^m$, if $f:U\rightarrow\mbb{C}^n$ satisfies for some $k$, $f(x_1,\cdots,x_k,\cdots,0)\equiv 0$, then there exists biholomorphic $\varphi:\mbb{C}^n\rightarrow\mbb{C}^n$ such that $\varphi\circ f=(0,\cdots,0,g_{k+1},\cdots,g_{n})$.}



%\bibliographystyle{IEEEtran}
%
%\bibliography{refs}
\end{document}

